%! Author = Kevin Lin

% Preamble
\documentclass[11pt,a4paper,margin=1in]{article}

% Packages
\usepackage{amsmath}
\usepackage{amssymb}
\usepackage{enumerate}
\usepackage{pdfpages}

\title{Notes}
\author{}
\date{}

% Document
\begin{document}
\maketitle

\section*{9/24/26}
Transform original event of taylor polynomials by applying a Markov inequality to
obtain an upper bound dependent on the moment generating function of our random variable.
By indepenence of the random variables, we can write the expectation of the exponential
of the sum as a product of each individual expectation. Then we can optimize this
final bound in terms of $\lambda$.
\\
\\
Bernoulli Sums:
\begin{gather*}
    X_1, \dots, X_n \text{ iid Bern}(p) \\
    E[X_i] = p, Var(X_i) = p(1 - p) \\
    P[X_i = 1] = p, P[X_i = 0] = 1 - p \\ 
    \mathcal{N}(0, 1) = \frac{1}{\sqrt{n}}\sum_{i = 1}^{n}\frac{X_i - E[X_i]}{\sqrt{Var(X_i)}} \\
    P[\frac{1}{n}\sum_{i=1}^{n}X_i - p \geq t] = P[\uparrow \geq t\sqrt{\frac{n}{p(1 - p)}}] \approx 1 - \Phi(t\sqrt{\frac{n}{p(1 - p)}}) \approx e^{-\frac{nt^2}{p(1 - p)}} \\
\end{gather*}

\noindent Large Deviation Principle:
\begin{gather*}
    P[\frac{1}{n}\sum_{i=1}^{n}X_i - p \geq t] \approx e^{-nD_{KL}(p + t || p)} \\
    D_{KL}(q || p) = E_q[\log(\frac{q}{p})] = (1 - q)\log(\frac{1 - q}{1 - p}) + q\log(\frac{q}{p}) \approx \frac{(q - p)^2}{2p(1 - p)}, q = t \\
\end{gather*}
by Chebyshev inequality:
\begin{align*}
    P[\frac{1}{n}\sum_{i=1}^{n}X_i - p \geq t] &\leq P[(\frac{1}{n}\sum_{i=1}^{n}X_i - p)^2 \geq t^2] \\
    &\leq \frac{1}{t^2}E[(\frac{1}{n}\sum_{i=1}^{n}X_i - p)^2] \\ 
    &= \frac{1}{t^2}\frac{p(1 - p)}{n} \\
\end{align*}
in the case of polynomial decay $k$:
\begin{align*}
    P[\frac{1}{n}\sum_{i=1}^{n}X_i - p \geq t] &\leq P[(\frac{1}{n}\sum_{i=1}^{n}X_i - p)^k \geq t^k] \\
    &\leq \frac{1}{t^k}E[(\frac{1}{n}\sum_{i=1}^{n}X_i - p)^k] \\ 
    &= \left(\frac{1}{t^2}\frac{p(1 - p)}{n}\right)^{\frac{k}{2}} \\
\end{align*}

Let $\psi(\lambda) = \log E[e^{\lambda(X_i - p)}]$ where it's MGF is $e^{\psi(\lambda)} = E[e^{\lambda(X_i - p)}]$, thus:
\begin{align*}
    P[\frac{1}{n}\sum_{i=1}^{n}X_i - p \geq t] &\leq P[e^{\lambda\sum_{i=1}^{n}(X_i - p)} \geq e^{\lambda nt}], \lambda > 0 \\
    &\leq e^{-\lambda nt}E[e^\lambda\sum_{i=1}^{n}(X_i - p)] \\
    &= e^{-\lambda nt}\prod_{i=1}^{n}E[e^\lambda(X_i - p)] \\
    &= e^{-\lambda nt + n\psi(\lambda)} \\
    &= e^{-n(\lambda t - \psi(\lambda))} \\
\end{align*}
Maximizing $\lambda t - \psi(\lambda)$ gives the best bound (minimizes due to $e^{-n(\dots)}$) \\
Given $\psi(\lambda) = \log E[e^{\lambda X_i}] - \lambda p = \log(1 - p + pe^\lambda) - \lambda p$:
\begin{align*}
    \lambda t - \psi(\lambda) &= \lambda t - \log(1 - p + pe^\lambda) + \lambda p \\
    &= \lambda(t + p) - \log(1 - p + pe^\lambda), \text{ minimize w/ derivative = 0}\\
    &= D_{KL}(p + t || p) \\
\end{align*}
Thus the bound is $P[\frac{1}{n}\sum_{i=1}^{n}X_i - p \geq t] \leq e^{-nD_{KL}(p + t || p)}$.
With this we can characterize how fast the probability decays as $n$ increases. \\

\noindent This gives us the Chernoff bound, where for i.i.d. $X_i, \dots, X_n$ where
$E[X_i] = 0, \psi(\lambda) = \log E[e^{\lambda X_i}]$:
\[
    \psi^*(t) = \sup_{\lambda > 0}(\lambda t - \psi(\lambda)) \\
\]
Then $P[\frac{1}{n}\sum_{i=1}^{n}X_i \geq t] \leq e^{-n\psi^*(t)}$.

\begin{gather*}
    X_i \sim \mathcal{N}(0, \sigma^2), \psi(\lambda) = \frac{\sigma^2\lambda^2}{2} \\
    \sup_{\lambda > 0} (\lambda t - \frac{\sigma^2\lambda^2}{2})
\end{gather*}

\section*{9/29/26}
Sub-Gaussian random variable: \\
Given a mean ($\mu$ = 0) random variable $X$ with variance $\sigma^2$ sub Gaussian,
if $E[e^{\lambda X}] \leq e^{\frac{\sigma^2\lambda^2}{2}}$, then 
$\psi(\lambda) \leq \frac{\lambda^2\sigma^2}{2}$. The random variable is consequently
bounded (tails are not heavier than a Gaussian distribution).

\begin{align*}
    P[|X| > t] &\leq 2 e^{-\frac{t^2}{2\sigma^2}}, \forall t \geq 0 \\
    P[X \geq t] &= P[e^{\lambda X} \geq e^{\lambda t}] \\
    &\leq e^{-\lambda t}E[e^{\lambda X}] \\
    &\leq e^{-\lambda t + \frac{\sigma^2\lambda^2}{2}} \\
    &= e^{-\frac{t^2}{2\sigma^2}} \\
\end{align*}
This is a case of exponential decay, standard for sub Gaussian random variables. \\

\noindent 
For $X_i$ i.i.d. $\sigma_i^2$ sub Gaussian, $E[e^{\lambda X_i}] \leq e^{\frac{\sigma_i^2\lambda^2}{2}}$,
then $\sum_{i=1}^{n}X_i$ is $(\sum_{i=1}^{n}\sigma_i^2)$ sub Gaussian, thus
\begin{gather*}
    P[|\sum_{i=1}^{n}X_i| > t] \leq 2e^{-\frac{t^2}{2\sum_{i=1}^{n}\sigma_i^2}} \\
\end{gather*}
As long as we can prove an upper bound of the MGF $E[e^{\lambda \sum_{i=1}^{n}X_i}] = E[\prod_{i=1}^{n}e^{\lambda X_i}] = \prod_{i=1}^{n}E[e^{\lambda X_i}]$.
We know that $E[e^{\lambda X_i}] \leq e^{\frac{\sigma_i^2\lambda^2}{2}}$, thus
\begin{gather*}
    E[e^{\lambda \sum_{i=1}^{n}X_i}] = \prod_{i=1}^{n}e^{\frac{\sigma_i^2\lambda^2}{2}} \leq e^{\frac{\lambda^2}{2}\sum_{i=1}^{n}\sigma_i^2} \\
\end{gather*}
This is a case of exponential decay, standard for sub Gaussian random variables. \\ 

\noindent
Special case of bounded random variables (Hoeffding's lemma): \\
If $a \leq X \leq b$ then $X - E[X]$ is $\frac{(b - a)^2}{4}$ sub Gaussian. \\
If $a_i \leq X_i \leq b_i$ and $X_i$ are independent, then $P[|\sum_{i=1}^{n}X_i - E[X_i]| \geq t] \leq 2e^{-\frac{2t^2}{\sum_{i=1}^{n}(b_i - a_i)^2}}$. \\
If a random variable is bounded, it's sub Gaussian is also always upper bounded. \\

\noindent
Suppose we have a Bernoulli random variable $X_1, \dots, X_n \sim \text{Bernoulli}(p)$ where $X_i \in \{0, 1\}$, 
then each $X_i$ is $\frac{1}{4}$ sub Gaussian. Using this result we can plug this
into (2) of the previous section to get a bound on the sum of Bernoulli random variables.
\begin{gather*}
    P[|\frac{1}{n}\sum_{i=1}^{n}X_i - p| \geq t] \leq 2e^{-\frac{2nt^2}{\sum_{i=1}^{n}(1 - 0)^2}} = 2e^{-2nt^2} \\
\end{gather*}
If $p = 0$, then each $\sum_{i=1}^{n}X_i = 0$. We can consequently define the sub 
Gaussian norm of a random variable $X$ as $\|X\|_{\psi_2} = \inf\{t > 0: E[e^{\frac{X^2}{t^2}}] \leq 2\}$, 
which is the smallest $t$ such that $E[e^{\frac{X^2}{t^2}}] \leq 2$. This is a 
measure of how "sub-Gaussian" a random variable is. \\

\noindent
\subsection*{Martingale method}
The concentration of $f(X_1, \dots, X_n)$ where $M_i = E[f(X_1, \dots, X_n) | X_1, \dots, X_i]$ 
is a martingale.
\begin{gather*}
    f(X_1, \dots, X_n) - E[f(X_1, \dots, X_n)] = M_n - M_0 = \sum_{i=1}^{n}(M_i - M_{i-1}) \\
\end{gather*}
This is a sum of Martingale differences. Then by Azuma-Hoeffding inequality, if
$X_1, \dots, X_n$ are independent then $E[e^{\lambda(M_i - M_{i-1})} | X_1, \dots, X_{i-1}] \leq e^{\frac{\lambda^2\sigma_i^2}{2}}$.
Then we can also show that $\sum_{i=1}^{n}(M_i - M_{i-1})$ is $(\sum_{i=1}^{n}\sigma_i^2)$ sub Gaussian, thus
$P[|f(X_1, \dots, X_n) - E[f(X_1, \dots, X_n)]| \geq t] \leq 2e^{-\frac{t^2}{2\sum_{i=1}^{n}\sigma_i^2}}$. \\

Proof:
\begin{align*}
    E[e^{\lambda(M_i - M_{i-1})}] &= E[E[e^{\lambda\sum{i=1}^{n}(M_i - M_{i-1})} | X_1, \dots, X_{i-1}]] \\
    &= E[e^{\lambda\sum_{i=1}^{n}(M_i - M_{i-1})}E[e^{\lambda(M_i - M_{i-1})} | X_1, \dots, X_{i-1}]] \\
    &\leq e^{\frac{\lambda^2\sigma_i^2}{2}}E[e^{\lambda\sum_{i=1}^{n}(M_i - M_{i-1})}] \\
\end{align*}

\noindent
\textbf{Corollary:} \\
Given $f(X_1, \dots, X_n)$, then $||\text{Diff}||_{\infty} = \sup_{x_1, \dots, x_n, x_i'}[\sup_{z} f(x_1, \dots, x_{i-1}, z, x_{i+1}, \dots, x_n) - \inf_{z} f(x_1, \dots, x_{i-1}, z, x_{i+1}, \dots, x_n)]$.
Then $P[|f(X_1, \dots, X_n) - E[f(X_1, \dots, X_n)]| \geq t] \leq 2e^{-\frac{2t^2}{\sum_{i=1}^{n}||\text{Diff}||_{\infty}^2}}$.
\end{document}